% Drop-in draft for the unfinished sections of the RTSS paper.
% Expected packages: graphicx (for \resizebox) and cite.

\begin{abstract}
Classifier cascades reduce average inference cost by accepting easy inputs at
small models and forwarding uncertain inputs to larger models.  Their usual
sequential execution, however, leaves processors idle and limits throughput
when inputs arrive continuously.  This paper presents a dynamic parallel IDK
(``I don't know'') cascade for image classification.  A lightweight model
handles every input first.  For an uncertain input, a learned router selects
one of two downstream models from the upstream model's confidence, entropy, and
top-two probability margin.  A scheduler coordinates model workers through an
asymmetric doubly linked queue, permits safe work stealing, and bounds the
number of in-flight inputs.  On ImageNet-V2, a Random-Forest skip decision
reduces cached sequential compute time by 5.9\% without changing top-1
accuracy at the reported precision.  In controlled cache replay, two in-flight
inputs increase throughput from 57.94 to 96.28 images/s while increasing mean
latency from 17.26 to 20.51 ms.  A live Apple-M1 implementation sustains a
30-frame/s input stream, whereas a 60-frame/s configuration is overloaded and
accumulates backlog.  These results show that learned routing and bounded
parallelism can improve average throughput, but also that admission rate and
shared-accelerator contention must be modeled explicitly for real-time use.
\end{abstract}

\begin{IEEEkeywords}
classifier cascades, real-time inference, dynamic scheduling, work stealing,
edge computing
\end{IEEEkeywords}

\section{Introduction}
\label{sec:introduction}

Modern image classifiers trade increasing execution cost for increasing
accuracy.  This trade-off is undesirable for continuous perception: many
frames are easy enough for a small model, but the hardest frames still require
a larger model.  An IDK classifier exposes this distinction by returning a
class only when its confidence exceeds a threshold and returning IDK
otherwise.  An IDK cascade runs increasingly capable classifiers until one
returns a class \cite{wang2018idk}.  It therefore avoids unnecessary work on
easy inputs without modifying the underlying pretrained models.

Sequential cascades save computation but expose little parallelism.  While a
large model processes one difficult input, the small model could already
process later inputs.  Simply assigning one model to each worker is also
insufficient: the data-dependent route creates unequal queues, accelerator
workers contend for a shared device, and unbounded admission converts a small
service-rate deficit into rapidly growing latency.  The scheduling problem is
therefore not only which classifier should run, but also when each input should
be admitted and which idle worker may execute it.

We study a dynamic parallel cascade built from pretrained ResNet-18,
ResNet-34, ResNet-50, and ResNet-152 classifiers.  The design combines
confidence-based early acceptance, a learned downstream router, bounded
in-flight execution, and asymmetric work stealing.  The paper makes three
contributions:

\begin{itemize}
    \item a low-overhead routing policy using confidence, entropy, and margin
    rather than confidence alone;
    \item a multiprocessing scheduler that preserves input--result association
    while allowing different models to process different inputs concurrently;
    and
    \item an evaluation that separates cached policy replay from live
    end-to-end execution and explicitly measures overload under fixed-rate
    camera arrivals.
\end{itemize}

The results demonstrate a throughput--latency trade-off rather than an
unqualified speedup.  Bounded concurrency substantially improves throughput;
unbounded concurrency achieves slightly higher throughput but produces
second-scale queueing delay.  The fixed-rate experiments further show that
completion throughput alone is insufficient: backlog slope and drain time are
needed to determine whether a real-time stream is sustainable.

\section{Background and Related Work}
\label{sec:background}

\subsection{IDK and Early-Exit Classification}

IDK cascades compose independently trained models and stop after the first
sufficiently confident answer \cite{wang2018idk}.  This is related to
early-exit networks such as BranchyNet \cite{teerapittayanon2016branchynet},
but the unit of composition differs.  Early-exit networks add intermediate
heads to one network and generally require joint training.  Our system instead
uses unmodified pretrained models, allowing model placement and scheduling to
change without retraining the classifiers.

Recent edge experiments likewise arrange pretrained classifiers from
lightweight to complex and use confidence thresholds to reduce average
inference time \cite{ratul2026accuracy}.  Our work differs by learning the
branch between downstream models and by evaluating queue stability under a
continuous periodic input stream.

Confidence is convenient but imperfect.  Modern neural networks may be
miscalibrated, so a large softmax probability does not necessarily equal the
true probability of correctness \cite{guo2017calibration}.  We therefore treat
the confidence threshold as an experimental policy parameter, report accuracy
separately from the non-IDK rate, and use entropy and margin as additional
router features.

\subsection{Scheduling Classifier Cascades}

Prior real-time analysis orders IDK classifiers to minimize expected duration,
optionally subject to a deadline \cite{baruah2023ordering}.  Later work models
arbitrary statistical dependence between classifiers and synthesizes static
multiprocessor list schedules \cite{abdelzaher2023scheduling}.  These models
provide an important offline baseline.  Our focus is complementary: a stream
of independent inputs is present simultaneously, observed service times vary,
and logically separate workers share one physical accelerator.  The scheduler
therefore makes online dispatch decisions and limits admission to control
queueing delay.

Work stealing is a standard way to move work from a busy worker to an idle one
\cite{blumofe1999workstealing}.  In this cascade, stealing is intentionally
asymmetric.  The larger downstream model may take a job assigned to the smaller
downstream model when its own side of the queue is empty.  The reverse steal is
disabled because a sample routed to the larger model is not assumed to be
classifiable by the smaller one.

\subsection{Limitations of Existing Approaches}

Three limitations motivate the proposed design.  First, a strictly sequential
cascade cannot overlap different inputs.  Second, a static multiprocessor
allocation does not react to route imbalance or runtime variation.  Third,
maximizing throughput without admission control can hide overload behind an
ever-growing queue.  Our design addresses these limitations operationally; it
does not claim a hard worst-case response-time guarantee.

\section{Models Under Comparison}
\label{sec:comparisons}

We compare five configurations.  All use ImageNet-pretrained residual networks
\cite{he2016resnet}; the exact model set and confidence threshold are reported
with every result.

\subsection{Sequential Learned-Skip Cascade}

The single-process baseline executes ResNet-18 first.  If it returns IDK, a
Random Forest predicts whether ResNet-34 should run or be skipped in favor of
the final heavy model.  All inference and routing operations execute
sequentially on one accelerator.

\subsection{Pipelined Sequential Cascade}

The pipelined baseline assigns ResNet-18, ResNet-34, ResNet-50, and ResNet-152
to separate logical workers while preserving the fixed cascade order.  A model
forwards an input only after returning IDK.  Different inputs may occupy
different stages concurrently, but an input cannot bypass a stage.

\subsection{Parallel Three-Model Cascades}

ResNet-18 processes every input.  A confident result terminates immediately;
otherwise, a learned router sends the input to ResNet-34 or to one heavy model.
We evaluate both ResNet-50 and ResNet-152 as the heavy alternative.  The two
downstream workers share a doubly linked waiting queue and the heavy worker may
steal a ResNet-34 job when safe.

\subsection{Parallel Four-Model Cascade}

ResNet-18 and ResNet-34 are started speculatively on separate logical workers.
If ResNet-18 succeeds, queued ResNet-34 work is canceled when possible and a
late result is ignored.  If both early models are uncertain, the router uses
ResNet-34's output features to choose ResNet-50 or ResNet-152.  This
configuration exposes more parallelism at the cost of speculative work and
greater contention.

\section{Experimental Setup}
\label{sec:setup}

\subsection{Dataset}

We use the three 10,000-image ImageNet-V2 test sets introduced by Recht
\emph{et al.} \cite{recht2019imagenetv2}.  Each contains ten images for each of
the 1,000 ImageNet classes.  The \emph{top-images} set contains the candidates
with the highest selection frequency, the \emph{threshold-0.7} set samples
images selected by at least 70\% of annotators, and the
\emph{matched-frequency} set matches the selection-frequency distribution of
the original ImageNet validation set.  The artifact keys confirm that the
three local sets are disjoint.  Router training uses matched-frequency and
top-images (20,000 images total); all reported policy and runtime tests use the
held-out threshold-0.7 set (10,000 images).

Images are resized to 256 pixels, center-cropped to $224\times224$, converted
to tensors, and normalized with the ImageNet statistics associated with the
pretrained weights.  Inference uses batch size one for the live experiments.

\subsection{Hardware and Software}

Experiments run on an Apple M1 MacBook Pro with four performance cores, four
efficiency cores, and 16 GB unified memory.  Classifiers execute through
PyTorch \cite{paszke2019pytorch} using the Metal Performance Shaders (MPS)
backend.  Each classifier has a separate Python process and queue, but all MPS
processes share the same physical accelerator; a ``worker'' therefore denotes
a logical execution context, not a dedicated GPU.  The inspected environment
uses Python 3.14, PyTorch 2.12, torchvision 0.27, NumPy 2.4, and scikit-learn
1.9.

\subsection{Router Training}

For probability vector $p$, the router feature vector is
\begin{equation}
  x(p) = \left[\max_i p_i,\; -\sum_i p_i\log(p_i+10^{-12}),\;
  p_{(1)}-p_{(2)}\right],
\end{equation}
where $p_{(1)}$ and $p_{(2)}$ are the largest and second-largest
probabilities.  Training includes only inputs for which the upstream model is
IDK and at least one downstream model reaches the confidence threshold.  When
both downstream models qualify, the smaller model is preferred.

The SVM uses an RBF kernel, $C=2$, automatic gamma scaling, and balanced class
weights \cite{cortes1995svm}.  Random Forest and Extra Trees use 100 trees,
maximum depth six, minimum leaf size 20, balanced class weights, and random
seed 42 \cite{breiman2001randomforests,geurts2006extratrees}.  The learned-skip
baseline instead uses 50 Random-Forest trees, maximum depth four, and minimum
leaf size 40.

\subsection{Metrics}

Top-1 accuracy is the fraction of all inputs whose final class equals the
ground-truth class.  Throughput is completed inputs divided by wall-clock run
time.  End-to-end latency is completion time minus admission time; we report
its mean and 95th percentile where recorded.  For fixed-rate input, we also
report the maximum number of unfinished frames, the backlog slope after a 10\%
warm-up interval, and drain time after the final arrival.  A run keeps up when
the realized arrival rate is within 5\% of its target and backlog growth is no
greater than $\max(0.5,0.02r)$ frames/s for target rate $r$.

\section{Results}
\label{sec:results}

\subsection{Learned Skip Decision}

Table~\ref{tab:skip} isolates the skip policy with cached logits and timing.
All three policies produce the same top-1 accuracy to three decimals.  The
Random-Forest policy skips ResNet-34 on 3,880 of the 6,278 inputs rejected by
ResNet-18 and reduces mean cached work from 18.33 to 17.25 ms per image, a
5.9\% reduction.  Its skip-decision accuracy is 69\%, compared with 38\% for a
fixed 0.3 skip threshold.

\begin{table}[t]
\caption{Cached skip-policy evaluation (R18/R34/R152, $\tau=0.9$, 10,000 images).}
\label{tab:skip}
\centering
\small
\begin{tabular}{lrrrr}
\hline
Policy & Accuracy & Work (ms/img) & B accepted & B skipped \\
\hline
No skip & 0.789 & 18.33 & 1,179 & 0 \\
Threshold & 0.789 & 17.80 & 1,130 & 1,277 \\
Random Forest & 0.789 & 17.25 & 807 & 3,880 \\
\hline
\end{tabular}
\end{table}

\subsection{Bounded Parallelism}

In controlled cached replay, the sequential RF baseline reaches 57.94
images/s with 17.26-ms mean latency.  Limiting the scheduler to two in-flight
inputs reaches 96.28 images/s (1.66$\times$) with 20.51-ms mean and 47.94-ms
95th-percentile latency.  Raising the limit to four reaches 104.84 images/s,
but mean latency increases to 37.64 ms.  With no limit, throughput saturates at
105.05 images/s while mean latency grows to 13.31 s and the 95th percentile to
57.55 s.  Thus, the useful operating point is bounded by latency, not by peak
throughput.

\subsection{Live Worker Results}

Table~\ref{tab:live} reports live preprocessing, routing, queueing, and MPS
inference.  These rows characterize implemented configurations; they are not
all causal comparisons because the model set and threshold vary.  The
three-model SVM configuration has the highest observed live throughput
(54.08 images/s) and low mean latency (27.60 ms), but also lower accuracy than
the four-model configuration.  The four-model SVM recovers 77.48\% accuracy at
42.21 images/s.  Its speculative ResNet-34 path cancels 1,900 jobs before
execution and ignores 1,682 completed or late jobs after ResNet-18 succeeds,
quantifying the cost of speculation.

\begin{table*}[t]
\caption{Live end-to-end results on threshold-0.7 ImageNet-V2.}
\label{tab:live}
\centering
\small
\begin{tabular}{llrrrrr}
\hline
Configuration & Models & $\tau$ & Accuracy (\%) & Throughput & Mean latency (ms) & Wall time (s) \\
\hline
Sequential RF skip & R18/R34/R50 & 0.9 & 77.49 & 25.87 & 36.56 & 386.53 \\
Pipelined sequential & R18/R34/R50/R152 & 0.7 & 76.84 & 36.30 & 8,606.36 & 275.51 \\
Parallel RF route & R18/R34/R152 & 0.7 & 77.18 & 41.15 & 3,339.89 & 243.02 \\
Parallel SVM route & R18/R34/R50 & 0.7 & 73.18 & 54.08 & 27.60 & 184.92 \\
Parallel four-model SVM & R18/R34/R50/R152 & 0.9 & 77.48 & 42.21 & 114.68 & 236.88 \\
\hline
\end{tabular}
\end{table*}

The very large mean latencies of the pipelined sequential and ResNet-152 RF
rows are queueing effects: their service rate falls below their effective
admission rate.  They should not be interpreted as single-image model
execution times.  This distinction explains why a run can show tens of
milliseconds of per-model execution yet seconds of end-to-end latency.

\subsection{Router Comparison}

Table~\ref{tab:router} evaluates all routers from the same cached training and
test features for the three-model 0.7-threshold configuration.  Router label
accuracy measures agreement with the confidence-derived routing target, not
ground-truth image correctness.  Random Forest best reproduces that target,
whereas Extra Trees gives the highest final cascade accuracy.  This mismatch
shows that a routing objective based only on whether a model crosses a
confidence threshold is not identical to the system objective.  Future router
training should optimize correctness and cost directly.

\begin{table}[t]
\caption{Cached router comparison ($\tau=0.7$).}
\label{tab:router}
\centering
\small
\begin{tabular}{lrrr}
\hline
Router & Route acc. (\%) & Cascade acc. (\%) & Work (ms/img) \\
\hline
SVM & 72.94 & 73.52 & 6.27 \\
Random Forest & 79.53 & 73.04 & 6.06 \\
Extra Trees & 70.65 & 73.88 & 6.35 \\
\hline
\end{tabular}
\end{table}

\subsection{Scheduler Behavior}

The in-flight limit is the main latency control.  At a limit of two, workers
overlap useful work without building a long queue.  Higher limits keep the
bottleneck worker busy but add waiting time almost linearly.  Work stealing is
beneficial only when the heavy worker is idle and the stolen sample was
originally assigned to the smaller downstream model.  No reverse stealing is
allowed.  Every result carries its input index, so out-of-order completion does
not change prediction--label association.

\section{Fixed-Rate Camera Experiment}
\label{sec:camera}

\subsection{Scenario and Method}

The offline workload always has another image available and therefore measures
maximum completion rate.  A camera instead produces periodic arrivals.  We
admit frames from an absolute monotonic clock, process every admitted frame,
and do not drop frames.  This experiment measures whether the queue remains
stable, not only how quickly the system drains after input ends.

\subsection{Results}

Table~\ref{tab:camera} shows two recorded operating points.  The 30-frame/s
R18/R34/R50 configuration keeps up: realized arrivals and completions both
equal approximately 30 frames/s, backlog has negative long-run slope, and only
17 ms is needed to drain after the last arrival.  Its transient maximum of 103
unfinished frames and 235-ms 95th-percentile latency show that keeping up on
average does not imply a per-frame 33.3-ms deadline.

The 60-frame/s R18/R34/R152 configuration does not keep up.  It completes 43.11
frames/s, has 1,776 unfinished frames at the final arrival, and requires 65.32
s to drain.  Mean latency reaches 19.13 s.  The two rows use different heavy
models and are therefore operating points rather than a controlled rate-only
comparison; nevertheless, the 60-frame/s result directly demonstrates
overload for that deployed configuration.

\begin{table*}[t]
\caption{Fixed-rate camera results ($\tau=0.7$, 10,000 frames; R18 and R34 are common to both rows).}
\label{tab:camera}
\centering
\small
\begin{tabular}{llrrrrrrrl}
\hline
Rate & Heavy & Acc. (\%) & In (fps) & Out (fps) & Mean (ms) & P95 (ms) & Backlog & Drain (s) & Stable \\
\hline
30 & R50 & 73.18 & 30.000 & 30.001 & 115.57 & 234.66 & 2 & 0.017 & yes \\
60 & R152 & 76.42 & 60.000 & 43.110 & 19,131.79 & 63,104.73 & 1,776 & 65.316 & no \\
\hline
\end{tabular}
\end{table*}

\section{Threats to Validity and Limitations}
\label{sec:limitations}

The experiments use one Apple M1 system, and multiple MPS processes contend
for one integrated GPU.  Results may differ on devices with dedicated
accelerators or true concurrent kernels.  Each configuration is represented by
one recorded run, so run-to-run variance and confidence intervals are not yet
available.  The study reports observed response times rather than WCET bounds
and consequently does not establish hard real-time schedulability.

Thresholds and heavy models differ across some configurations.  We report
these parameters explicitly and avoid attributing their combined effect to the
scheduler alone.  Router training and testing use disjoint ImageNet-V2
variants, but all variants follow the same collection process; evaluation on
additional datasets is needed.  Finally, no frame-dropping or backpressure
policy is applied.  The overload experiment therefore measures queue growth
under lossless processing, not the behavior of a production camera pipeline
that may discard stale frames.

\section{Conclusion and Future Work}
\label{sec:conclusion}

This paper presented a dynamic parallel IDK cascade that combines learned
routing, logical model workers, a partitioned doubly linked queue, asymmetric
work stealing, and bounded admission.  The experiments support three
conclusions.  First, a learned skip policy can remove intermediate-model work
without reducing measured accuracy.  Second, limited overlap yields a large
throughput gain, whereas unbounded overlap converts bottleneck utilization into
unacceptable queueing delay.  Third, a camera-rate claim must be supported by
backlog stability: the tested system sustains 30 frames/s but not the recorded
60-frame/s configuration.

Future work will calibrate model confidence on a held-out calibration split,
train routers against a joint correctness--latency objective, and repeat every
configuration across multiple devices and seeds.  A practical real-time
version should add deadline-aware admission, stale-frame dropping, and explicit
backpressure.  Profiling shared-accelerator interference and incorporating it
into response-time analysis are also necessary steps toward hard real-time
guarantees.
